%==============================================================
\section{Biên dịch như một bài toán quyết định tuần tự}
\label{sec:mdp}
%==============================================================

Thuật toán~\ref{alg:routing} cho thấy routing chỉ có \textbf{đúng một} chỗ cần quyết định, lặp lại
nhiều lần: chọn cạnh nào để chèn SWAP. Cấu trúc này khớp tự nhiên với một Markov decision process
(MDP) \cite{sutton2018rl} -- không phải vì ``RL giải được mọi thứ'', mà vì bản chất bài toán đúng là
như vậy. Chương này phát biểu MDP đó, trình bày các quyết định thiết kế khi hiện thực nó thành một
môi trường \code{gymnasium}, và đặt thang mốc so sánh.

\subsection{Phát biểu MDP cho qubit routing}

\begin{table}[h]
  \centering
  \caption{Các thành phần MDP của bài toán qubit routing.}
  \label{tab:mdp}
  \small
  \begin{tabularx}{\linewidth}{@{}lX@{}}
    \toprule
    \textbf{Thành phần} & \textbf{Trong routing} \\
    \midrule
    State $s_t$ & front layer $F$ của DAG còn lại, mapping hiện tại $\pi$, (và lớp lookahead) \\
    Action $a_t$ & chọn một cạnh $(a,b)\in E$ của coupling map để chèn SWAP; $|\mathcal{A}| = |E|$ \\
    Transition & chèn SWAP $\to$ cập nhật $\pi$ $\to$ \textbf{bắt buộc} thực thi mọi cổng đã hợp lệ (tất định) \\
    Reward $r_t$ & $-1$ mỗi SWAP, hoặc $-\Delta\text{depth}$, hoặc $\log$ fidelity của SWAP \\
    Terminal & DAG rỗng (\code{terminated}); vượt số bước tối đa (\code{truncated}) \\
    \bottomrule
  \end{tabularx}
\end{table}

\paragraph{Tính Markov.} Cặp $(F,\pi)$ cùng DAG còn lại xác định đầy đủ tương lai -- không cần biết đã
đi đường nào để tới trạng thái hiện tại. Đây là điều kiện bắt buộc để áp dụng RL, và nó được thoả.

\paragraph{Thực thi cổng là bắt buộc, không phải action.} Nếu để agent tự quyết ``thực thi hay hoãn''
một cổng đã hợp lệ, action space lớn hơn mà không thêm khả năng nào: hoãn một cổng đã hợp lệ không bao
giờ có lợi. Vì vậy transition tự động thực thi mọi cổng chạy được, và agent chỉ quyết định SWAP.

\subsection{Thiết kế observation}

Cách encode state quyết định agent có học được hay không. Môi trường \code{RoutingEnv} dùng ba quyết
định thiết kế:
\begin{enumerate}[label=(\alph*), noitemsep]
  \item \textbf{Dùng qubit vật lý, không dùng qubit logic.} Observation ghi ``hai qubit \emph{vật lý}
    $Q_a, Q_b$ đang cần kề nhau'' dưới dạng ma trận kề $n\times n$. Nhờ vậy observation \emph{bất biến}
    với cách đánh số qubit logic; nếu không, agent phải học $n!$ phiên bản của cùng một tình huống.
  \item \textbf{Thêm một lớp lookahead.} Chỉ nhìn front layer là đủ tính Markov nhưng khó học: một SWAP
    tốt cho hiện tại có thể tệ cho các cổng ngay sau. Heuristic \code{lookahead} hơn \code{basic}
    chính vì lý do này~\eqref{eq:hlook}.
  \item \textbf{Thêm một số đo tiến độ} -- tỉ lệ cổng đã thực thi -- để agent biết mình ở đầu hay cuối
    episode.
\end{enumerate}
Observation vì vậy là vector $[\,\text{front } (n\times n) \mid \text{lookahead } (n\times n) \mid
\text{tiến độ}\,]$, với \code{line(5)} có kích thước $2\cdot5^2+1 = 51$.

\paragraph{Action masking.} Ở mỗi bước chỉ một số cạnh là có nghĩa -- SWAP trên cạnh không liên quan
tới qubit nào đang bị kẹt chắc chắn chỉ làm mạch dài thêm. Môi trường trả mask qua
\code{info["action\_mask"]} nhưng vẫn nhận action bị mask (mọi thư viện RL cần \code{step()} không
raise). Trên \code{line(5)}, trung bình 69\% action là có nghĩa nên mask ít tác dụng; nhưng với
heavy-hex 115 qubit và 132 cạnh, mask là bắt buộc.

\subsection{Thiết kế reward}

Reward định nghĩa cái agent sẽ giỏi, nên đây là quyết định quan trọng nhất (Bảng~\ref{tab:reward}).

\begin{table}[h]
  \centering
  \caption{Ba kiểu reward trong \code{RoutingEnv}. $\epsilon_{ab}$ là sai số cổng hai qubit trên
  cạnh $(a,b)$; hệ số 3 vì một SWAP bằng 3 cổng hai qubit.}
  \label{tab:reward}
  \small
  \begin{tabularx}{\linewidth}{@{}llXX@{}}
    \toprule
    \textbf{Kiểu} & \textbf{Công thức} & \textbf{Ưu điểm} & \textbf{Nhược điểm} \\
    \midrule
    \code{swap} & $-1$ mỗi SWAP & đơn giản, dễ so với các công trình khác & bỏ qua depth và noise \\
    \code{depth} & $-\Delta\text{depth}$ & thưởng SWAP song song được & phụ thuộc thứ tự thực thi \\
    \code{fidelity} & $3\log(1-\epsilon_{ab})$ & gần mục tiêu thật nhất; học được cạnh nào tốt & cần dữ liệu sai số, phụ thuộc thiết bị \\
    \bottomrule
  \end{tabularx}
\end{table}

Ba reward này \textbf{không cùng chiều}: thí nghiệm ở Mục~\ref{sec:exp-topology} cho thấy
\code{grid(3,3)} cần ít SWAP hơn \code{ring(9)} nhưng depth lại cao hơn. Vì vậy ``agent giảm 10\% số
SWAP'' không tự động có nghĩa là mạch chạy tốt hơn trên thiết bị thật. Reward \code{fidelity} là log
của xác suất thành công~\eqref{eq:psuccess}, nên tổng reward của một episode chính là
$\log P_{\text{thành công}}$ phần do routing gây ra.

\subsection{Cổng kiểm tra: hợp lệ \emph{và} đúng}

Trong môi trường RL, hai lỗi ở Mục~\ref{sec:verify} còn nguy hiểm hơn, vì agent \emph{tối ưu hoá
đúng cái ta đo}: nếu môi trường có lỗ hổng, agent sẽ khai thác đúng lỗ hổng đó. Hàm
\code{kiem\_tra\_ket\_qua} kiểm tra cả hai điều kiện cho mỗi episode: (1) \textbf{hợp lệ} bằng
\code{CheckMap}; (2) \textbf{đúng} bằng so sánh unitary với mạch gốc sau khi bù hoán vị bằng
$\pi^{-1}$. Chạy 100 episode cho mỗi policy mốc, cả \code{RandomPolicy} lẫn \code{GreedyPolicy} đều
đạt 100\% hợp lệ và 100\% đúng. Môi trường cũng qua \code{gymnasium.utils.env\_checker.check\_env}
(đúng API 5 giá trị của \code{step()}, \code{reset(seed=...)} tái lập được, observation đúng
\code{observation\_space}).

\subsection{Hai cách phát biểu khác}

Routing không phải bài toán biên dịch duy nhất phát biểu được thành bài toán quyết định
(Bảng~\ref{tab:framings}).

\begin{table}[h]
  \centering
  \caption{Ba cách phát biểu compilation thành bài toán học tăng cường.}
  \label{tab:framings}
  \small
  \begin{tabularx}{\linewidth}{@{}lXXX@{}}
    \toprule
    & \textbf{Routing} & \textbf{Synthesis (T-count)} & \textbf{Pass ordering} \\
    \midrule
    State & $(F, \pi)$ + lookahead & unitary còn lại / mạch Clifford+T hiện tại & metric của mạch (depth, size, số CX) \\
    Action & chọn cạnh để SWAP & thêm một cổng / áp một luật rewrite & chọn pass tiếp theo \\
    Reward & $-1$, $-\Delta$depth, $\log$ fidelity & $-1$ mỗi cổng $T$, thưởng khi đạt $\epsilon$ & $-\Delta$(hàm cost) \\
    Baseline & \code{SabreSwap(decay)} & \code{gridsynth}; \code{OptimizeCliffordT} ($-34\%$) & \code{optimization\_level} 0$\to$3: 37$\to$19 CX \\
    Độ khó & trung bình; action rời rạc & khó: state liên tục, reward thưa & dễ nhất; action space nhỏ \\
    \bottomrule
  \end{tabularx}
\end{table}

Synthesis là hướng của AlphaTensor-Quantum \cite{ruiz2025alphatensor} và RL cho compiling một qubit
\cite{moro2021rlcompiling}: khó hơn vì không gian state liên tục và reward thưa, nhưng giá trị thực tế
lớn vì T-count chi phối tài nguyên của máy chịu lỗi. Pass ordering có action space nhỏ nhất (vài chục
pass) và episode ngắn, nên là điểm khởi đầu tốt để kiểm tra toàn bộ pipeline thí nghiệm.
